
\section{Built in functions}

Some Built in functions in GCC(remember that if u are using ll, need to add ll at the end):
\begin{center}
\begin{tabular}{ c | l }
  \_\_builtin\_popcount(x)  & Counts the quantity of one’s(set bits) in an integer.  \\ 
 \_\_builtin\_parity(x) & Checks the Parity of a integer.\\
 
 &Returns true(1) if odd parity(odd quantity of set bits)\\
 
 & Returns false(0) for even parity(even quantity of set bits).  \\  
 \_\_builtin\_clz(x) &  Counts the leading quantity of zeros of the integer.\\
 \_\_builtin\_ctz(x) & Counts the trailing quantity of zeros of the integer.  \\ 
 \_\_builtin\_popcountll(x)  & Counts the quantity of one’s(set bits) in an long long.  \\ 
 \_\_builtin\_parityll(x) & Checks the Parity of a long long.\\
 
 &Returns true(1) if odd parity(odd quantity of set bits)\\
 
 & Returns false(0) for even parity(even quantity of set bits).  \\   
 \_\_builtin\_clzll(x) &  Counts the leading quantity of zeros of the long long.\\
 \_\_builtin\_ctzll(x) & Counts the trailing quantity of zeros of the long long.  \\ 
\end{tabular}
\end{center}
clz can be used to find the first set bit, since u can use:
\begin{lstlisting}
int x;
ll xl;
int fsi =  31 - __builtin_clz(x);
int fsll = 63 - __builtin_clzll(xl);
\end{lstlisting}
